\section{Conclusion and Future Work}
Feature implementation is a critical frontier in autonomous software engineering. However, existing feature-level benchmarks suffer from two primary limitations: unrealistic task inputs, and data leakage risks. Typically relying on technical implementation hints (e.g., function signatures) or static datasets with narrow scope, these benchmarks fail to align with the realistic challenges and workflows inherent in real-world software development scenarios.
To address these challenges, we introduce FeatBench, a novel benchmark designed to evaluate feature implementation based on natural language requirements that are strictly devoid of code hints. FeatBench utilizes a fully automated, evolving pipeline to construct tasks from the latest open-source repositories, ensuring the benchmark remains dynamically updated and aligned with the broad spectrum of real-world development.
A comprehensive evaluation of 157 tasks across 27 diverse repositories demonstrates that implementing features from natural language poses a substantial engineering challenge, with SOTA agents achieving a resolved rate of only 29.94\%. Furthermore, our analysis uncovers a prevalence of ``scope creep,'' where agents exhibit an aggressive implementation behavior. While this behavior can occasionally yield superior software architecture, it paradoxically introduces critical regressions and over-engineering by diverging from the user's explicit intent.

In the future, we will expand FeatBench's scope beyond Python to include other popular programming languages, such as Java and Go. We also plan to explore mechanisms to align agent behaviors with realistic engineering practices, aiming to balance strict specification adherence with the capability for robust software design.